\begin{claim}{E.19}{Known}{A coherent replacement for arbitrary relative path composition}
Let $\mathcal M$ be $\U$-small and let $W$ be a subcategory containing every object. The hammock localization, followed on each mapping simplicial set by $\operatorname{Ex}^{\infty}$, gives $\U$-small Kan complexes
\[
 \mathcal H_W(a,b)=\operatorname{Ex}^{\infty}
          L^H(\mathcal M,W)(a,b)
 \qquad(a,b\in\mathcal M)
\]
with strictly associative, strictly unital simplicial composition. Their component category is canonically $\mathcal M[W^{-1}]$. Precomposition or postcomposition by any specified arrow in $W$ is a weak equivalence of these mapping complexes. These are coherent mapping objects for the relative localization. No comparison with a specified ordinary colimit of the manuscript's path categories is part of this assertion.
\end{claim}
\begin{proof}
The hammock simplices, their faces, and their composition are those of \citet[Definition~2.1]{DwyerKan1980Calculating}: a simplex is a finite commuting hammock, its leftward and vertical arrows lie in $W$, and composition concatenates horizontally and reduces consecutive same-direction columns and identity columns. The resulting simplicial category has strictly associative composition and identity hammocks. Every set of finite hammocks lies in $\U$, and taking the countable union of lengths and simplicial degrees preserves that size. Proposition~3.1 of the same paper identifies its component category with the ordinary localization. Proposition~3.3 proves the weak-equivalence assertion for composition with an arrow of $W$.

The functor $\operatorname{Ex}$ preserves products because it is a right adjoint; filtered colimits of sets commute with finite products, so its sequential iteration $\operatorname{Ex}^{\infty}$ also preserves finite products. Applying it to the composition maps therefore gives
\[
 \mathcal H_W(b,c)\times\mathcal H_W(a,b)
       \longrightarrow\mathcal H_W(a,c)
\]
with the same strict associativity and unit equations. The natural maps from a simplicial set to its $\operatorname{Ex}^{\infty}$ replacement are weak equivalences and the replacements are Kan. Thus the component category and the specified transport equivalences are unchanged. In a model-category presentation these are the derived function complexes by \citet[Proposition~4.4 and Corollary~4.7]{DwyerKan1980Function}. This supplies an unconditional relative-category construction in the stated small universe, while comparison with any other raw path model remains a separate mathematical assertion, as recorded in E.26--E.27.
\end{proof}
